% TODO: Related Work
%
% Structure into three subsections:
%  1. CXL simulation frameworks (CXLMemSim, CXL-DMSim, ESF) — note ESF
%     already shows adaptive-vs-oblivious routing helps under a
%     noisy-neighbor scenario; differentiate this paper as mechanism-level
%     characterization of the saturation threshold plus a lightweight,
%     Python-based rapid-prototyping tool.
%  2. Datacenter congestion-aware load balancing (CONGA, Hedera,
%     queue-occupancy ECMP variants) — establish occupancy-weighted ECMP
%     lineage; note it hasn't been evaluated specifically in a CXL fabric
%     context.
%  3. CXL multi-tenancy/fairness (Equilibria) — note this operates at the
%     OS/memory-tiering layer, complementary to (not overlapping with)
%     this paper's fabric-routing-layer focus if H2 is included.

\section{Related Work}
\label{sec:related-work}

\subsection{CXL Simulation Frameworks}
TODO: discuss CXLMemSim~\cite{yang2023cxlmemsim}, CXL-DMSim
\cite{wang2025cxldmsim}, and ESF~\cite{an2024esf}. Note that ESF already
shows adaptive-vs-oblivious routing helps under a noisy-neighbor scenario;
this paper differentiates itself via mechanism-level characterization of
the saturation threshold and a lightweight, Python-based rapid-prototyping
tool (TraceCXL).

\subsection{Congestion-Aware Load Balancing}
TODO: discuss CONGA~\cite{alizadeh2014conga} and
Hedera~\cite{alfares2010hedera}, plus queue-occupancy-based ECMP
variants~\cite{ye2018wecmp} (note: this is a P4/weighted-ECMP paper, not a
paper specifically titled ``QALL''/``QLLB'' — see the
\texttt{\% TODO: verify} comment above this entry in
\texttt{references.bib} for what's still unconfirmed). Establish the
occupancy-weighted ECMP lineage and note it has not previously been
evaluated in a CXL fabric context.

\subsection{CXL Multi-Tenancy and Fairness}
TODO: discuss Equilibria~\cite{zhao2026equilibria}. Note this operates at
the OS/memory-tiering layer, complementary to (not overlapping with) this
paper's fabric-routing-layer focus, if H2 (cross-tenant fairness, see
experiment spec) is included in the final paper.
